\documentclass[a4paper,11pt,fleqn]{scrartcl}
\usepackage{../../Styles/Exercice}

\newcommand{\chnum}{Chapitre 8}
\newcommand{\chtitle}{Réactions chimiques forcées}
\newcommand{\dtype}{Corrigé}

\begin{document}
\maketitle
\begin{multicols}{2}

\section{Schématiser les transferts d'électrons aux électrodes}

\textbf{1. Demi-équations électroniques}

On observe :
\begin{itemize}
    \item Formation de dihydrogène à l’électrode B $\Rightarrow$ réduction des ions $\ce{H+}$
    \item Formation de dioxygène à l’électrode A $\Rightarrow$ oxydation de l’eau
\end{itemize}

\textbf{À la cathode (réduction, électrode B)} :
\[
\ce{2H+(aq) + 2e^- -> H2(g)}
\]

\textbf{À l’anode (oxydation, électrode A)} :
\[
\ce{2H2O(l) -> O2(g) + 4H+(aq) + 4e^-}
\]

\textbf{2. Sens des électrons et bornes}

\begin{itemize}
    \item Les électrons circulent dans le circuit extérieur de l’anode vers la cathode.
    \item L’anode est reliée au pôle positif du générateur.
    \item La cathode est reliée au pôle négatif du générateur.
\end{itemize}

Donc :
\begin{itemize}
    \item Électrode A : anode (+)
    \item Électrode B : cathode (-)
\end{itemize}

\section{Déterminer les variations de quantité de matière}

\textbf{1. Quantité d’électricité}

\[
Q = I \times \Delta t
\]

Avec :
\[
\Delta t = 20\,\text{min} = 1200\,\text{s}
\]

\[
Q = 0{,}40 \times 1200 = 480\,\text{C}
\]

\textbf{2. Quantité de matière d’électrons}

\[
n(e^-) = \frac{Q}{\mathcal{F}} = \frac{480}{96500}
\]

\[
n(e^-) \approx 4{,}97 \times 10^{-3}\,\text{mol}
\]

\textbf{3. Quantités de matière formées}

\textbf{Pour le rhodium :}
\[
\ce{Rh^{3+} + 3e^- -> Rh}
\]

\[
n(\ce{Rh}) = \frac{n(e^-)}{3} \approx \frac{4{,}97 \times 10^{-3}}{3}
\]

\[
n(\ce{Rh}) \approx 1{,}66 \times 10^{-3}\,\text{mol}
\]

\textbf{Pour le dioxygène :}
\[
\ce{2H2O -> O2 + 4H+ + 4e^-}
\]

\[
n(\ce{O2}) = \frac{n(e^-)}{4} \approx \frac{4{,}97 \times 10^{-3}}{4}
\]

\[
n(\ce{O2}) \approx 1{,}24 \times 10^{-3}\,\text{mol}
\]

\textbf{4. Masse et volume}

\textbf{Masse de rhodium :}
\[
m = n \times M = 1{,}66 \times 10^{-3} \times 102{,}9
\]

\[
m \approx 0{,}171\,\text{g}
\]

\textbf{Volume de dioxygène :}
\[
V = n \times V_m = 1{,}24 \times 10^{-3} \times 24{,}0
\]

\[
V \approx 2{,}98 \times 10^{-2}\,\text{L} \approx 29{,}8\,\text{mL}
\]

\section{Nickelage électrolytique d'une pièce métallique}

\textbf{1. Électrode de la pièce}

La pièce doit être reliée au pôle négatif (cathode), car il y a réduction :

\[
\ce{Ni^{2+}(aq) + 2e^- -> Ni(s)}
\]

\textbf{2. Nature de l’électrode}

La pièce est la \textbf{cathode}.

\textbf{3. Rôle de l’électrode en nickel}

L’électrode en nickel joue le rôle d’anode soluble :

\[
\ce{Ni(s) -> Ni^{2+}(aq) + 2e^-}
\]

Cela permet de maintenir constante la concentration en ions $\ce{Ni^{2+}}$.

\textbf{4. Quantités de matière}

\[
n(\ce{Ni}) = \frac{m}{M} = \frac{3{,}0}{59{,}0}
\]

\[
n(\ce{Ni}) \approx 5{,}08 \times 10^{-2}\,\text{mol}
\]

D’après la demi-équation :

\[
n(e^-) = 2 \times n(\ce{Ni}) \approx 1{,}02 \times 10^{-1}\,\text{mol}
\]

\textbf{5. Durée de l’électrolyse}

\[
Q = n(e^-) \times \mathcal{F}
\]

\[
Q = 1{,}02 \times 10^{-1} \times 96500 \approx 9{,}84 \times 10^{3}\,\text{C}
\]

\[
\Delta t = \frac{Q}{I} = \frac{9{,}84 \times 10^{3}}{8{,}0}
\]

\[
\Delta t \approx 1230\,\text{s} \approx 20{,}5\,\text{min}
\]

\end{multicols}
\end{document}